% !PRINTED 534
\weberpage{534}

%<<CONT
die $a_{1,r}, a_{2,r}, \ldots, a_{r-1,r}$ können alle gleich Null
angenommen werden.

Bezeichnen wir mit $\Delta$ die Grundzahl des Körpers $\Omega$, so
ergiebt sich die Discriminante des durch (3) bestimmten Systemes
$\alpha_1, \alpha_2, \ldots, \alpha_n$ nach der Formel \S.\,144, (11):
\weq{5}{
\Delta(\alpha_1, \alpha_2, \ldots, \alpha_n) = a_{1,1}^2\, a_{2,2}^2
\ldots a_{n,n}^2\, \Delta.
}
Dies ist eine von Null verschiedene ganze rationale Zahl, und folglich
sind die Grössen $\alpha$ eine Basis von $\Omega$.

Um also zu zeigen, dass das so bestimmte System $\alpha_1, \alpha_2,
\ldots, \alpha_n$ eine Basis von $\mu$ ist, bleibt noch nachzuweisen,
dass jede durch $\mu$ theilbare ganze Zahl $\alpha$ in der Form (2)
dargestellt werden kann. Nehmen wir, um diesen Beweis zu führen, irgend
einen Index $r \leqq n$ an, und suchen die Bedingung dafür, dass eine
ganze Zahl von der Form
\weq{6}{
\gamma_r = h_1\,\omega_1 + h_2\,\omega_2 + \cdots + h_r\,\omega_r
}
durch $\mu$ theilbar ist, wenn $h_1, h_2, \ldots, h_r$ ganze rationale
Zahlen sind.

Zunächst folgt, dass $h_r$ durch $a_{r,r}$ theilbar sein muss. Denn
bezeichnen wir mit $q_r$ den Rest der Division von $h_r$ durch $a_{r,r}$,
und setzen
\[
h_r = l_r\,a_{r,r} + q_r, \qquad 0 \leqq q_r < a_{r,r},
\]
so ergiebt sich nach (6)
\[
\gamma_r - l_r\,\alpha_r = (h_1 - l_r\,a_{1,r})\,\omega_1 +
(h_2 - l_r\,a_{2,r})\,\omega_2 + \cdots + q_r\,\omega_r,
\]
und diese Zahl müsste auch durch $\mu$ theilbar sein. Dies ist aber nach
der Definition von $a_{r,r}$ nur möglich, wenn $q_r = 0$ ist. Dann aber
erhält $\gamma_r - l_r\,\alpha_r$ den Ausdruck:
\weq{7}{
\gamma_r - l_r\,\alpha_r = h_1'\,\omega_1 + h_1'\,\omega_2 + \cdots +
h_{r-1}'\,\omega_{r-1}, % [sic]
}
wird also von derselben Form, wie (6), nur dass $r - 1$ an Stelle von $r$
tritt, und in (7) ist dieselbe Schlussweise zu wiederholen. Demnach
ergiebt sich durch vollständige Induction der Satz:

\gesperrt{Eine in der Form $h_1\,\omega_1 + h_2\,\omega_2 + \cdots +
h_r\,\omega_r$ darstellbare ganze Zahl des Körpers $\Omega$ ist immer
dann und nur dann durch $\mu$ theilbar, wenn sie in der Form}
\[
l_1\,\alpha_1 + l_2\,\alpha_2 + \cdots + l_r\,\alpha_r
\]
\gesperrt{darstellbar ist, in der die Co\"efficienten $l_1, l_2, \ldots,
l_r$ ganze rationale Zahlen sind.}

Setzt man in diesem Satze $r = n$, so hat man den Beweis dafür, dass das
Zahlensystem (3) eine Basis von $\mu$ ist.
